\begin{trapbox}
	\begin{description}[style=nextline, leftmargin=2.8em, labelsep=0.5em, itemsep=0.35em]
		\item[\textbf{\textcolor{CTrap}{易错点一（底数忽略）}}] 求导 $(-2)^x$ 或 $0^x$ 时错误套用 $a^x$ 公式.
		\item[\textbf{\textcolor{CTrap}{正确理解（底数范围）：}}] 公式 $(a^x)' = a^x\ln a$ 要求 $a>0$ 且 $a\neq 1$. 当 $a\le 0$ 或 $a=1$ 时，该公式不适用.
		\item[\textbf{\textcolor{CTrap}{错因分析（条件遗忘）：}}] 只记公式形式，忽略公式的成立前提.
	\end{description}

	\medskip
	\begin{description}[style=nextline, leftmargin=2.8em, labelsep=0.5em, itemsep=0.35em]
		\item[\textbf{\textcolor{CTrap}{易错点二（角度误解）}}] 认为 $\sin x^\circ$ 的导数仍然是 $\cos x^\circ$.
		\item[\textbf{\textcolor{CTrap}{正确理解（弧度要求）：}}] 导数公式 $(\sin x)'=\cos x$, $(\cos x)'=-\sin x$ 仅在 $x$ 以弧度为单位时成立. 若 $x$ 以度为单位，导数需乘以 $\frac{\pi}{180}$.
		\item[\textbf{\textcolor{CTrap}{错因分析（单位忽略）：}}] 默认角度制下导数公式不变，忽略了弧度制的必要性.
	\end{description}

	\medskip
	\begin{description}[style=nextline, leftmargin=2.8em, labelsep=0.5em, itemsep=0.35em]
		\item[\textbf{\textcolor{CTrap}{易错点三（符号混淆）}}] 记忆余弦导数为 $\sin x$，丢掉负号.
		\item[\textbf{\textcolor{CTrap}{正确理解（余弦导数）：}}] $(\cos x)' = -\sin x$，务必保留负号.
		\item[\textbf{\textcolor{CTrap}{错因分析（形似记混）：}}] 正弦和余弦导数形式相似，符号容易混淆.
	\end{description}
\end{trapbox}
